% T1_1TDZ.pdf -- Cvičení 1: 1. termodynamický zákon
\providecommand{\dbar}{{\mathchar'26\mkern-12mu d}}

\section{Cvičení: 1.\ TDZ}
\strana{T1-1}

\begin{itemize}
  \item \kw{rovnovážný stav} \dots\ makroskopicky statický, charakterizovaný malým počtem TD parametrů \\
  \underline{Př:} jednoduchý systém: \quad $p, V, N$ \\
  \hspace*{2cm} 0.TDZ \quad $\Rightarrow$ \quad $\vartheta$ \\
  \hspace*{2cm} 1.TDZ \quad $\Rightarrow$ \quad $U$ \qquad $(\Delta U = W^{ad} = W + Q)$
  \item ne $\forall$ nezávislé (svázány stav.\ rovnicemi) $\Rightarrow$ \underline{stavový prostor} \\
  \underline{Př:} uzavřený jednoduchý systém ($N = $ konst) \\
  \hspace*{2cm} $\underline{p, V}$ \quad \& \quad $U = U(p,V)$ \quad \& \quad $\vartheta = \vartheta(p,V)$
  \item \kw{kvazistat.\ proces} \dots\ posloupnost infinit.\ blízkých rovnovážných stavů
  \[
    \underset{\textcolor{NavyBlue}{\text{úplný dif.}}}{\dd U} = \underset{\textcolor{OliveGreen}{\text{neúplné dif.}}}{\dbar W + \dbar Q}
  \]
  \underline{Př:} \quad $\dbar W = -p\,\dd V$ \hfill
  $\begin{aligned}
    \dbar W < 0 &\ \dots\ \text{systém práci koná} \\
    \dbar W > 0 &\ \dots\ \text{systém práci přijímá}
  \end{aligned}$
  \item \underline{klasifikace dějů} \hfill
  \begin{tabular}[t]{lll}
    isochorický & \dots & $V = $ konst $\Leftrightarrow \dd V = 0$ \\
    isobarický & \dots & $p = $ konst $\Leftrightarrow \dd p = 0$ \\
    isotermický & \dots & $\vartheta = $ konst $\Leftrightarrow \dd\vartheta = 0$ \\
    adiabatický & \dots & $\dbar Q = 0$
  \end{tabular}
\end{itemize}

\begin{priklad}[Příklad 1.*]
rovnice adiabaty pro plyn popsaný stavovou rovnicí $U = cpV + \alpha$
\end{priklad}

\begin{priklad}[Příklad 3.*]
Určete práci a teplo podél libovolné trajektorie mezi dvěma body ve stavovém prostoru, ležícími na téže adiabatě. Rovnice adiabaty systému je $p^3V^5 = \text{konst}$. [Callen str.\ 21, cv.\ 1]
\end{priklad}

\begin{priklad}[1) $k = 2$]
$\dd\omega = (x^2 + y)\,\dd x - x\,\dd y = 0$
\begin{enumerate}[label=\alph*)]
  \item ukažte, že lze hledat int.\ faktor ve tvaru $\mu = \mu(x)$
  \item ukažte, že $\mu = \dfrac{1}{x^2}$
  \item ukažte $\dd\sigma = \mu\,\dd\omega \ \Rightarrow\ \sigma = x + \dfrac{y}{x^2} + C$
\end{enumerate}
\end{priklad}

\begin{priklad}[2) $k = 3$]
$\dd\omega = yz\,\dd x + xz\,\dd y + xyz\,\dd z$
\begin{enumerate}[label=\alph*)]
  \item ukažte, že $\exists\,\mu$ \quad (rce 4)
  \item ověřte $\mu = \dfrac{1}{xyz}$
  \item ukažte $\sigma = \log x + \log y + z$
\end{enumerate}
\end{priklad}

%----------------------------------------------------------------------
\strana{T1-2}
\subsection*{1) TD stav (rovnovážný)}
\begin{itemize}
  \item \underline{ext.\ veličiny}/parametry \qquad $U/S, V, N, P, M, \dots$
  \item \underline{int.\ \quad --\,{}"\,--} \qquad\qquad\quad\ $T, p, \mu, \mathcal{E}, H, \dots$
  \item \underline{úplný popis} \dots\ úplná sada $k$ ext.\ parametrů (typ.\ $U, V, N_i$) \\
  \hspace*{2.2cm} \dots\ (později uvidíme) lze téměř lib.\ jinou kombinací $k$ parametrů
  \begin{itemize}[nosep]
    \item[--] $k$ \dots\ \# TD st.\ volnosti
    \item[--] alesp.\ 1 param.\ musí být extenzivní
  \end{itemize}
  \item stav TD \underline{rovnováhy} \dots\ tranzitivní (0.TDZ) \\
  \quad \dots\ stac.\ stav (neměnnost parametrů) bez stac.\ proudů
  \pozn{šipka: odkaz z bodu 1) na tento bod}
  \item \underline{kvazistat.\ proces} \dots\ efektivně \uv{nekonečně} pomalý vzhledem k mikroskop.\ čas.\ škálám (ie, relax.\ časům) \\
  \quad $\Rightarrow$ lze popisovat jako \underline{posloupnost rovn.\ stavů} \\
  \quad (\pozor{NB: v TD popisu obvykle nevystupuje čas!})
\end{itemize}

\subsection*{2) TD systém}
\begin{itemize}
  \item popsán \underline{stav.\ rovnicemi} -- závislost vnitřních stav.\ proměnných ($U, p, \mu, P, M$) na vnějších a $T$
  \begin{itemize}[nosep]
    \item \underline{vnější} -- externě kontrolovatelné
  \end{itemize}
  \item kalor.\ stav.\ rce \quad $U = U(T, V, N, \dots)$ \\
        termická \hspace*{1.55cm} $p = p(T, V, N, \dots)$, \ $\mu = \mu(T, V, N, \dots)$, \dots
\end{itemize}

\subsection*{3) práce}
změna vnitřní energie \underline{spojená} se (kontrolovatelnou) \underline{změnou vnějších parametrů}
\begin{enumerate}[label=\alph*)]
  \item mech: \quad $W = -p\,\dd V$
  \item chem: \quad $W = \mu\,\dd N$ \quad $(= \sum\mu_i\,\dd N_i)$
  \item elstat: \quad $W = \mathcal{E}\,\dd P$
  \item mg: \quad $W = H\,\dd M$
\end{enumerate}

%----------------------------------------------------------------------
\strana{T1-3}
\begin{veta}[1.\ TDZ]
\[ \vzorec{\dd U = \dbar Q + \dbar W} \]
\begin{itemize}[nosep]
  \item znam.\ notace: \quad $> 0 \Rightarrow$ systém energii přijal \\
        \hspace*{2.75cm} $< 0 \Rightarrow$ \quad --\,{}"\,-- \quad odevzdal
\end{itemize}
\end{veta}

$U$ \dots\ stav.\ veličina, \ $\dd U$ úplný diferenciál \\
$Q, W$ \dots\ dějové veličiny, \ $\dbar Q, \dbar W$ neúplné dif.

\pozor{!} \quad $A\to B$: \quad $\Delta U$, \ ale $Q_{AB}, W_{AB}$

$\Rightarrow$ \kw{operační definice $Q$}: \quad
$\dbar Q = \underset{\substack{\uparrow\\ \text{umím \uv{určit}}\\ \text{(použitím vhodných stěn)}}}{\dd U} - \underset{\substack{\uparrow\\ \text{umím měřit}\\ \text{(spojeno se změnou vnějš.\ parametru)}}}{\dbar W}$

\begin{itemize}
  \item \underline{konkrétně:} umím měřit práci $\Rightarrow$ umím určit $\Delta U$ v adiabat.\ procesu, $U$ urč.\ až na adit.\ konstantu
  $\Rightarrow$ zvolím ref.\ \sout{bod} stav $A$ s $U_A = 0$ a v principu znám $U$ všech stavů, které lze spojit s $A$ adiabat.\ procesem (\& nejsou $\forall$ stavy \dots\ $\nexists$ -- potřebujeme ještě např.\ $\Delta U$ v isochor.\ procesu \dots)
\end{itemize}

\noindent\rule{\linewidth}{0.4pt}

\begin{itemize}[label=$\Rightarrow$]
  \item Úloha č.\ 1 -- Callen p.\ 21, exercise 1
  \item adiabata pro \quad $U = \frac52 pV + \alpha$ \quad $\leftarrow$ kalorická stav.\ rovnice
  \item $\Delta$ cyklus pro \quad $\downarrow$
  \item integrace 2D formy / 3D formy?
\end{itemize}

%----------------------------------------------------------------------
\strana{T1-4}
\subsection*{Stavové rovnice}
\begin{itemize}[nosep]
  \item[--] obecně rovnice typu \quad $f_\alpha(T, V, N_i, \mu_i, p, \dots) = 0$ \\
  $=$ závislosti vnitřních stav.\ proměnných na $T$ a vnějších parametrech
  \begin{itemize}[nosep]
    \item vnější -- externě kontrolovatelné \dots\ $V, N_i, E, H, \dots$
    \item vnitřní \hspace*{4.2cm} $U, p, \mu_i, P, M, \dots$ \\
    \pozn{$\Rightarrow$ tato je spíše fenomenologická definice, obecně jsou stav.\ rce parc.\ derivace fund.\ rovnice (TD potenciálů)}
  \end{itemize}
\end{itemize}
\begin{tabular}{@{}ll}
kalorická rovnice: & $U = U(T, V, N)$ \\
termické \quad\ : & $p = p(\dots)$, \ $\mu_i = \mu_i(\dots)$
\end{tabular}

\begin{itemize}
  \item id.\ plyn: \quad $pV = Nk_BT$ \hfill $\bullet$ vdW \quad $\left(p + \dfrac{a}{v^2}\right)(v - b) = k_BT$ \\
  \hspace*{1.6cm} $U = c_VNk_BT$ \hfill $u = ?$\hspace*{2.2cm}
  \item \kw{koeficienty lin.\ odezvy}: derivace stav.\ parametrů (stav.\ rcí) \\
  \quad $\Leftrightarrow$ druhé derivace fundamentální rovnice
\end{itemize}
\begin{align*}
  \alpha &= \frac1V\left(\pdv{V}{T}\right)_p && \text{koef.\ tepl.\ roztažnosti} \\
  \beta &= -\frac1V\left(\pdv{V}{p}\right)_T \equiv \kappa_T && \text{izoterm.\ stlačitelnost} \\
  \gamma &= \frac1p\left(\pdv{p}{T}\right)_V && \text{koef.\ teplotní tuhosti}
\end{align*}
\[
  \Rightarrow\quad \vzorec{\alpha = p\beta\gamma} \qquad \text{(to teď nechci dělat)}
\]

\pozor{!} \quad $\vzorec{C_X = T\left(\pdv{S}{T}\right)_X = \left(\frac{\dbar Q}{\dd T}\right)_X}$ \\
\hspace*{4.3cm} $\updownarrow$ \\
Toto je jen zkrácený zápis \quad $\dfrac{\partial S(T,X)}{\partial T}$
\begin{itemize}[nosep]
  \item $S(T,X)$ \uv{$=$} $S(T,Y)$ \quad z fyz.\ hlediska -- jedná se o tutéž veličinu
  \item z mat.\ hlediska je to jiná funkce $\Rightarrow$ $\left(\pdv{S}{T}\right)_X \neq \left(\pdv{S}{T}\right)_Y$ \\
        ale $S(T,X) = S(T, Y(X,T))$
\end{itemize}
\pozn{vzhůru nohama v rohu: 0--5 \quad 5 nejlepší}

%----------------------------------------------------------------------
\strana{T1-5}
\begin{priklad}[Adiabata (pro id.\ plyn) -- \pozor{není stav.\ rovnice!}]
Mějme kalor.\ stav.\ rci $U = \frac52 pV + \alpha$. Najděte rovnici adiabaty (křivka v $p$-$V$, $V$-$T$, \dots\ diagramu, odpovídající ději $\dbar Q = 0$)
\end{priklad}
\[
  \text{1.TDZ}\ \Rightarrow\ \dbar Q = \dd U - \dbar W = 0 = \underbrace{\tfrac52 p\,\dd V + \tfrac52 V\,\dd p}_{\dd U} + p\,\dd V
\]
\begin{align*}
  &\Rightarrow\ \tfrac72 p\,\dd V + \tfrac52 V\,\dd p = 0 \\
  &\Rightarrow\ \tfrac75\frac{\dd V}{V} = -\frac{\dd p}{p}
   \ \Rightarrow\ \tfrac75\log V = -\log p + A \\
  &\Rightarrow\ \log\big(V^{7/5}p\big) = A \ \Rightarrow\ \vzorec{V^{7/5}p = \text{konst}} \\
  &\Leftrightarrow\ pV^\kappa = \text{konst} \qquad \kappa = \tfrac75 = \frac{c_p}{c_V}
\end{align*}
(je to 2-atom.\ id.\ plyn)

%----------------------------------------------------------------------
\strana{T1-6}
\begin{priklad}[Adiabata vs.\ izoterma]
Určete podmínky, za kterých adiabata splývá s izotermou (pro jednoduchý systém, charakt.\ $p, V, T$ (resp.\ 2 z nich))
\end{priklad}

\begin{minipage}{0.6\textwidth}
\underline{adiabata}: \quad $\dbar Q = 0 = \dd U + p\,\dd V$
\end{minipage}
\begin{minipage}{0.35\textwidth}
\fbox{$\begin{aligned} \&\ U &= U(T,V) \\ p &= p(T,V)\end{aligned}$}
\end{minipage}
\begin{align*}
  \Rightarrow\ 0 &= \left(\pdv{U}{T}\right)_V\dd T + \left[\left(\pdv{U}{V}\right)_T + p\right]\dd V \\
  &= C_V\,\dd T + \left[\left(\pdv{U}{V}\right)_T + p\right]\dd V
\end{align*}

\underline{izoterma}: \quad $\dd T = 0$, $\dd V \neq 0$ \\
\quad $\Rightarrow$ aby to platilo zároveň, musí být ($C_V\neq 0$)
\[
  \vzorec{\left(\pdv{U}{V}\right)_T + p = 0} \qquad \dots\ \text{rce pro } \left(\pdv{U}{V}\right)_T \quad (p = p(V,T))
\]
\underline{NB:} nelze splnit pro id.\ plyn \dots\ tam je $\left(\pdv{U}{V}\right)_T \equiv 0$ jinde než pro $T = 0$

\underline{vdW}:
\[
  \left(\pdv{u}{v}\right)_T = \frac{a}{v^2}
  \ \Rightarrow\ \frac{a}{v^2} + \frac{RT}{v-b} - \frac{a}{v^2} = 0
  \ \Rightarrow\ \frac{RT}{v-b} = 0 \ \Leftrightarrow\ T = 0
\]

\begin{veta}[3.\ zákon]
Za absolutní nuly \sout{izoentropické} jsou izotermické procesy zároveň adiabatické
\[ \lim_{T\to 0}\left(\pdv{S}{X}\right)_T = 0 \]
\end{veta}
\underline{id.\ plyn} \quad $\left(\pdv{u}{v}\right)_T = 0 \ \Rightarrow\ p = 0 = \dfrac{NRT}{V} \ \Rightarrow$ opět OK

%----------------------------------------------------------------------
\strana{T1-7}
\begin{priklad}[$\left(\pdv{T}{V}\right)_p\left(\pdv{V}{p}\right)_T\left(\pdv{p}{T}\right)_V = -1$]
(viz teorie implic.\ fcí)
\end{priklad}
\begin{itemize}
  \item uvažujme stav.\ rovnici \quad $f(T, V, p) = \text{konst}$ \quad \sout{$\Leftrightarrow p = p(T,V)$}
  \[
    \Rightarrow\ \dd f = \left(\pdv{f}{T}\right)\dd T + \left(\pdv{f}{V}\right)\dd V + \left(\pdv{f}{p}\right)\dd p = 0
  \]
\end{itemize}
1) $\Leftrightarrow p = p(T,V)$ \quad \& izobara \quad $\Rightarrow T = T(V) \Rightarrow \dd T = \frac{\dd T}{\dd V}\dd V$ \quad \sout{$T = T(V)$}
\begin{align*}
  &\Rightarrow\ \left(\pdv{f}{T}\right)\dd T + \left(\pdv{f}{V}\right)\dd V = 0 \\
  &\Rightarrow\ \left(\pdv{f}{T}\right)_{V,p}\left(\pdv{T}{V}\right)_p + \left(\pdv{f}{V}\right)_{T,p} = 0 \\
  &\Rightarrow\ \vzorec{\left(\pdv{T}{V}\right)_p = -\frac{\left(\pdv{f}{V}\right)_{T,p}}{\left(\pdv{f}{T}\right)_{V,p}}}
\end{align*}
\pozn{v originále je ve jmenovateli napsáno $(\partial f/\partial V)_{V,p}$ -- zřejmě překlep za $(\partial f/\partial T)_{V,p}$}

2) $T = T(p,V)$ \quad \& \quad $\dd T = 0$ \\
3) $V = V(T,p)$ \quad \& \quad $\dd V = 0$
\[
  \Rightarrow\ \vzorec{\left(\pdv{T}{V}\right)_p\left(\pdv{V}{p}\right)_T\left(\pdv{p}{T}\right)_V = -1}
\]

%----------------------------------------------------------------------
\strana{T1-8}
\textbf{1.TDZ:} \quad $\dd U = \dbar Q + \dbar W$ \qquad
$\dbar W = \underset{\text{mech}}{-p\,\dd V} \ \big/\ \underset{\text{el}}{+\mathcal{E}\,\dd P} \ \big/\ \underset{\text{mg}}{+H\,\dd M}$

\begin{itemize}
  \item umím měřit práci $\Rightarrow$ umím měřit $\Delta U$ v adiabat.\ procesu $\Rightarrow$ umím měřit $Q$ mezi týmiž dvěma stavy v neadiabat.\ procesu ($U$ stavová veličina, $W/Q$ \uv{dějové})
  \item adiabat.\ izol.\ systém: stav nelze změnit jinak, než změnou vnějších parametrů
\end{itemize}

\begin{priklad}[Př.\ č.\ 1]
\begin{itemize}[nosep]
  \item válec s pístem, ve kterém je uzavřen 1-komp.\ plyn
  \item experimentem zjistím (systém uzavřu adiabat.\ stěnami), že v adiabat.\ procesu platí
  \[ p^3V^5 = \text{konst.} \qquad (\text{pro } Q = 0) \text{ -- odpovídá 1-atom.\ ideálnímu plynu} \]
\end{itemize}
\end{priklad}

\begin{center}
\begin{tikzpicture}[>=Stealth,scale=1.1]
  \draw[->] (-0.3,0) -- (6,0) node[right] {$V$};
  \draw[->] (0,-0.3) -- (0,3.6) node[left] {$p$};
  \coordinate (A) at (1.2,3); \coordinate (D) at (5,3);
  \coordinate (C) at (1.2,0.5); \coordinate (B) at (5,0.5);
  \draw (0,3) node[left] {$p_A$} -- (0.1,3);
  \draw (0,0.5) node[left] {$p_B$} -- (0.1,0.5);
  \draw (1.2,0) node[below] {$V_A$} -- (1.2,0.1);
  \draw (5,0) node[below] {$V_B$} -- (5,0.1);
  \node[above left] at (A) {$A$}; \node[right] at (D) {$D$};
  \node[left] at (C) {$C$}; \node[right] at (B) {$B$};
  \draw[thick] (A) -- (D) node[midway,above] {$(Q>0)$};
  \draw[thick] (D) -- (B) node[midway,right] {$(Q<0)$};
  \draw[->,thick] (5,2) -- (5,1.6);
  \draw[thick] (A) -- (C);
  \draw[->,thick] (1.2,2) -- (1.2,1.6);
  \draw[thick] (C) -- (B);
  \draw[->,thick] (2.9,0.5) -- (3.3,0.5);
  \draw[->,thick] (2.9,3) -- (3.3,3);
  % přímka AB
  \draw[thick] (A) -- (B);
  \draw[->,thick] (3.0,1.82) -- (3.2,1.71);
  % adiabata
  \draw[thick] (A) .. controls (2,1.6) and (3.3,0.9) .. (B) node[pos=0.45,below left] {\small $Q=0$};
  \draw[->,thick] (2.6,1.4) -- (2.8,1.3);
  \node[align=center] at (8,2.2) {$A, B$ leží na\\ 1 adiabatě};
\end{tikzpicture}
\end{center}

\begin{enumerate}[label=\alph*)]
  \item (kvazi-statické) práce a teplo podél trajektorií \sout{AB}, \fbox{ADB}, ACB, \fbox{AB (přímka)} \\
  \hspace*{4cm} $\hookrightarrow p = p_A + \dfrac{V - V_A}{V_B - V_A}(p_B - p_A) \Rightarrow W = \frac12(p_A + p_B)(V_A - V_B)$
\end{enumerate}
\begin{enumerate}[label=\alph*),start=2]
  \item $\Delta U$ ($V = $ konst, $p\nearrow$) -- isochor.\ proces
  \item umíme již projít $\forall$ stavy?
\end{enumerate}

%----------------------------------------------------------------------
\strana{T1-9}
\textbf{a) I)} \quad $\Delta U_{AB}$ -- zjistíme z adiabat.\ procesu \hfill
\begin{minipage}[t]{0.33\textwidth}\small
\underline{NB:} $A, B$ musí ležet na adiabatě (netriv.\ požadavek)
\end{minipage}
\begin{align*}
  p_AV_A^{5/3} &= p_BV_B^{5/3} = pV^{5/3} \\
  \Rightarrow\ p &= p_A\left(\frac{V_A}{V}\right)^{5/3} \\
  W_{AB}^{ad} &= -\int_{V_A}^{V_B} p_A\left(\frac{V_A}{V}\right)^{5/3}\dd V = \frac32 p_AV_A^{5/3}\Big[V^{-2/3}\Big]_A^B \\
  &= \frac32 p_AV_A^{5/3}\Big[V_B^{-2/3} - V_A^{-2/3}\Big] = \Delta U_{AB} < 0
\end{align*}

\textbf{II)} \quad \underline{ADB:} \quad $\Delta U = \Delta U_{AB} < 0$
\begin{align*}
  W_{ADB} &= W_{AD} = p_A(V_A - V_B) < 0 \\
  \Delta U_{AB} &= W_{ADB} + Q_{ADB}
\end{align*}
\begin{minipage}{0.22\textwidth}\small
\pozor{!} neznám $Q_{AD}$ a $Q_{DB}$ samostatně
\end{minipage}
\begin{minipage}{0.76\textwidth}
\[
  \Rightarrow\ \vzorec{Q_{ADB} = \tfrac32 p_AV_A^{5/3}\Big[V_B^{-2/3} - V_A^{-2/3}\Big] + p_A(V_B - V_A)}
\]
\end{minipage}

\underline{NB:}
\begin{itemize}[nosep]
  \item výsledek jen zdánlivě nezávisí na $p_B$, neboť $p_B$ je \sout{dáno} určeno $V_B$ skrze adiabatu
  \item umíme říct, jestli se teplo dodá nebo odevzdá? (zde $Q > 0$, ale jak je to s tou plochou v $pV$-diagramu u cyklu? OK -- uzavřu cyklus, plyn koná práci při větším tlaku, je stlačován při menším $\boxtimes$)
\end{itemize}

%----------------------------------------------------------------------
\strana{T1-10}
\textbf{b)} alias \uv{co se stavy, které neleží na adiabatě} \\
Do válce umístíme lopatku (vrtulku), kterou umíme nějak bezkontaktně (mg) pohánět, přičemž známe moment síly (tzv.\ kroutivý moment) motorku $M$.
\begin{itemize}
  \item \underline{naměříme:} \quad $V = $ konst \ $\Rightarrow$ \ $\dfrac{\dd p}{\dd t} = \dfrac23\dfrac{\omega}{V}M$
\end{itemize}
$\Rightarrow$ ukažte, že takto umíme změřit $\Delta U$ pro lib.\ 2 stavy při stejném objemu \\
\quad (avšak pouze v procesu s rost.\ tlakem -- proč?)

\rule{3cm}{0.4pt}

\underline{řeš:} \quad výkon motorku: \quad $\dfrac{\dd U}{\dd t} = M\omega > 0$

\begin{minipage}{0.06\textwidth}\Huge\pozor{!}\end{minipage}
\begin{minipage}{0.9\textwidth}
(předanou en.\ nelze klasif.\ jako práci nebo teplo -- z hlediska plynu -- nejedná se o kvazistac.\ proces)
\end{minipage}
\begin{align*}
  &\Rightarrow\ \frac{\dd p}{\dd t} = \frac{2}{3V}\frac{\dd U}{\dd t} \ \Rightarrow\ \dd U = \frac32 V\,\dd p \geq 0 \\
  &\Rightarrow\ U_A - U_C = \frac32 V_A(p_A - p_B)
\end{align*}

\textbf{c)} libovolné 2 stavy lze vždy spojit adiabatou a isochorou (z nichž jedna je ve správném směru) $\Rightarrow$ umíme určit $\Delta U$ \sout{mezi} mezi lib.\ 2 stavy

\begin{minipage}[t]{0.3\textwidth}\small
\fbox{\parbox{0.9\linewidth}{\underline{NB:} může to být možné jen 1 směrem s adiabatou mezi stavy}}
\end{minipage}
\begin{minipage}[t]{0.65\textwidth}
(\underline{NB:} Joule -- právě toto experimentálně potvrdil)
\end{minipage}

%----------------------------------------------------------------------
\strana{T1-11}
\begin{priklad}[Př.\ 2]
Mějme \sout{systém} plyn, pro který platí $U = \frac52 pV + \alpha$, uzavřený ve stejném pístu a projedeme cyklus
\end{priklad}

\begin{center}
\begin{tikzpicture}[>=Stealth,scale=1.1]
  \draw[->] (-0.3,0) -- (6.2,0) node[right] {$V$};
  \draw[->] (0,-0.3) -- (0,4) node[left] {$p$};
  \coordinate (A) at (1,0.6); \coordinate (B) at (5.2,0.6); \coordinate (C) at (1,3.4);
  \draw (0,3.4) node[left] {$p_C$} -- (0.1,3.4);
  \draw (0,0.6) node[left] {$p_A$} -- (0.1,0.6);
  \draw (1,0) node[below] {$V_A$} -- (1,0.1);
  \draw (5.2,0) node[below] {$V_B$} -- (5.2,0.1);
  \node[below left] at (A) {$A$}; \node[below right] at (B) {$B$}; \node[above] at (C) {$C$};
  \draw[thick] (A) -- (B) -- (C) -- cycle;
  \draw[->,thick] (2.9,0.6) -- (3.2,0.6);
  \draw[->,thick] (3.4,1.8) -- (3.1,2.0);
  \draw[->,thick] (1,2.2) -- (1,1.9);
  \node[align=left] at (8.3,2.6) {$\sum W, Q$ na $\forall$\\ úsecích};
  \node at (8.3,1) {$\Rightarrow$ účinnost?};
\end{tikzpicture}
\end{center}

\underline{AB:}
\[
  \left.\begin{aligned}
    \Delta U_{AB} &= \tfrac52 p_A(V_B - V_A) > 0 \\
    W_{AB} &= -p_A(V_B - V_A) < 0
  \end{aligned}\right\}
  \ \Rightarrow\ Q_{AB} = \tfrac72 p_A(V_B - V_A) > 0
\]

\underline{BC:}
\begin{align*}
  &\vzorec{\Delta U_{BC} = \tfrac52(p_CV_A - p_AV_B)} \gtrless 0 \\
  W_{BC} &= -\int_{V_B}^{V_A} p\,\dd V = \Big/\, p = p_A + \frac{(p_C - p_A)(V - V_B)}{V_A - V_B} \,\Big/ \\
  &= p_A(V_B - V_A) - \frac12\frac{p_C - p_A}{V_A - V_B}\Big[(V_A - V_B)^2 - (V_B - V_B)^2\Big] \\
  &= p_A(V_B - V_A) + \frac12(p_C - p_A)(V_B - V_A) \\
  &\vzorec{W_{BC} = \tfrac12(p_A + p_C)(V_B - V_A)} > 0
\end{align*}
\[
  \Rightarrow\ \vzorec{Q_{BC} = \Delta U_{BC} - W_{BC}} \gtrless 0 = \tfrac12\big(p_AV_A - 6p_AV_B - p_CV_B + 6p_CV_A\big)
\]
\hfill \fbox{$Q_{BC} > 0 \Leftrightarrow (p_A + 6p_C)V_A > (6p_A + p_C)V_B$}

%----------------------------------------------------------------------
\strana{T1-12}
\begin{minipage}{0.62\textwidth}
\underline{CA:}
\begin{align*}
  \Delta U_{CA} &= \tfrac52 V_A(p_A - p_C) < 0 \\
  W_{CA} &= 0 \\
  \Rightarrow Q_{CA} &= \tfrac52 V_A(p_A - p_C) < 0
\end{align*}
\underline{celkem:} \quad $W > 0$, \ $\Delta U = 0$, \ $Q < 0$
\end{minipage}
\begin{minipage}{0.36\textwidth}
\pozor{!} \quad $W_{tot} = \frac12(V_A - V_B)(p_A - p_C)$ \\
\quad $\Rightarrow$ obsah $\Delta$!
\end{minipage}

$\Rightarrow$ můžeme doplnit znaménko u BC? \\ ačkoli na AB a CA neumíme jednoduše porovnat. \quad \sout{$Q_{AB} = \frac72 p_A$}

\begin{priklad}[Př.\ 3]
adiabata $p = p(V)$ pro daný systém
\end{priklad}
\begin{align*}
  \dd U &= \tfrac52 p\,\dd V + \tfrac52 V\,\dd p = -p\,\dd V \qquad \Leftarrow (\dbar Q = 0) \\
  &\Rightarrow\ \tfrac72 p\,\dd V = -\tfrac52 V\,\dd p \\
  &\Rightarrow\ 7\frac{\dd V}{V} = -5\frac{\dd p}{p} \ \Rightarrow\ 7\log V = -5\log p + K \\
  &\Rightarrow\ V^7 = \tilde Kp^{-5} \ \Rightarrow\ \vzorec{pV^{7/5} = \text{konst}} \\
  &\Rightarrow\ \text{2-atom.\ id.\ plyn}
\end{align*}
(\underline{NB:} $pV = Nk_BT \Rightarrow U = \frac52 Nk_BT$)

